\begin{ex}
Tính: $2 \times 1$.
\shortans{2}
\loigiai{
Ta có $2 \times 1 = 2$.
}
\end{ex}

\begin{ex}
Tính: $2 \times 2$.
\shortans{4}
\loigiai{
Ta có $2 \times 2 = 4$.
}
\end{ex}

\begin{ex}
Tính: $2 \times 3$.
\shortans{6}
\loigiai{
Ta có $2 \times 3 = 6$.
}
\end{ex}

\begin{ex}
Tính: $2 \times 4$.
\shortans{8}
\loigiai{
Ta có $2 \times 4 = 8$.
}
\end{ex}

\begin{ex}
Tính: $2 \times 5$.
\shortans{10}
\loigiai{
Ta có $2 \times 5 = 10$.
}
\end{ex}

\begin{ex}
Tính: $2 \times 6$.
\shortans{12}
\loigiai{
Ta có $2 \times 6 = 12$.
}
\end{ex}

\begin{ex}
Tính: $2 \times 7$.
\shortans{14}
\loigiai{
Ta có $2 \times 7 = 14$.
}
\end{ex}

\begin{ex}
Tính: $2 \times 8$.
\shortans{16}
\loigiai{
Ta có $2 \times 8 = 16$.
}
\end{ex}

\begin{ex}
Tính: $2 \times 9$.
\shortans{18}
\loigiai{
Ta có $2 \times 9 = 18$.
}
\end{ex}

\begin{ex}
Tính: $2 \times 10$.
\shortans{20}
\loigiai{
Ta có $2 \times 10 = 20$.
}
\end{ex}

\begin{ex}
Tính giá trị biểu thức: $2 \times 4 + 15$.
\shortans{23}
\loigiai{
Ta có $2 \times 4 = 8$, rồi $8 + 15 = 23$.
}
\end{ex}

\begin{ex}
Tính giá trị biểu thức: $2 \times 7 + 26$.
\shortans{40}
\loigiai{
Ta có $2 \times 7 = 14$, rồi $14 + 26 = 40$.
}
\end{ex}

\begin{ex}
Tính giá trị biểu thức: $2 \times 9 - 9$.
\shortans{9}
\loigiai{
Ta có $2 \times 9 = 18$, rồi $18 - 9 = 9$.
}
\end{ex}

\begin{ex}
Tính giá trị biểu thức: $2 \times 6 + 18$.
\shortans{30}
\loigiai{
Ta có $2 \times 6 = 12$, rồi $12 + 18 = 30$.
}
\end{ex}

\begin{ex}
Tính giá trị biểu thức: $2 \times 8 - 10$.
\shortans{6}
\loigiai{
Ta có $2 \times 8 = 16$, rồi $16 - 10 = 6$.
}
\end{ex}

\begin{ex}
Điền số thích hợp: $2 \times \dots = 16$.
\shortans{8}
\loigiai{
Số cần điền là: $16 : 2 = 8$.
}
\end{ex}

\begin{ex}
Điền số thích hợp: $2 \times \dots = 14$.
\shortans{7}
\loigiai{
Số cần điền là: $14 : 2 = 7$.
}
\end{ex}

\begin{ex}
Một đôi đũa có 2 chiếc đũa. Hỏi 7 đôi đũa có tất cả bao nhiêu chiếc đũa?
\shortans{14}
\loigiai{
Có: $2 \times 7 = 14$ chiếc đũa.
}
\end{ex}

\begin{ex}
Một đôi chim bồ câu có 2 con. Hỏi 9 đôi chim bồ câu có tất cả bao nhiêu con?
\shortans{18}
\loigiai{
Có: $2 \times 9 = 18$ con chim.
}
\end{ex}

\begin{ex}
Mỗi bạn học sinh được phát 2 quyển vở. Hỏi 10 bạn học sinh được phát bao nhiêu quyển vở?
\shortans{20}
\loigiai{
Phát: $2 \times 10 = 20$ quyển vở.
}
\end{ex}

\begin{ex}
Một con chim có 2 cái cánh. Hỏi 8 con chim có tất cả bao nhiêu cái cánh?
\shortans{16}
\loigiai{
Có: $2 \times 8 = 16$ cái cánh.
}
\end{ex}

\begin{ex}
Mỗi chiếc xe đạp cần 2 bánh xe. Hỏi để lắp ráp 6 chiếc xe đạp cần bao nhiêu bánh xe?
\shortans{12}
\loigiai{
Cần: $2 \times 6 = 12$ bánh xe.
}
\end{ex}

\begin{ex}
Trong vườn có 5 luống hoa, mỗi luống có 2 cây hoa hồng. Có tất cả bao nhiêu cây hoa hồng?
\shortans{10}
\loigiai{
Có: $2 \times 5 = 10$ cây hoa.
}
\end{ex}

\begin{ex}
So sánh: Tích $2 \times 8$ lớn hơn tích $2 \times 6$ bao nhiêu đơn vị?
\shortans{4}
\loigiai{
Ta có $16 - 12 = 4$.
}
\end{ex}

\begin{ex}
Điền số thích hợp vào ô trống: $2 \times 5 + 2 = 2 \times \dots$
\shortans{6}
\loigiai{
Vế trái là $10 + 2 = 12$, vậy $2 \times 6 = 12$. Số cần điền là 6.
}
\end{ex}
